% Generated draft: requires mathematical and source review.
\documentclass[11pt]{article}
\usepackage[margin=25mm]{geometry}
\usepackage{fontspec}
\setmainfont{Latin Modern Roman}
\usepackage{amsmath,amssymb,mathtools,unicode-math}
\setmathfont{Latin Modern Math}
\usepackage{enumitem,needspace,xurl,hyperref}
\setlength{\parindent}{0pt}
\setlength{\parskip}{5pt}
\setlength{\emergencystretch}{3em}
\newcommand{\N}{\mathbb{N}}
\newcommand{\Z}{\mathbb{Z}}
\newcommand{\Q}{\mathbb{Q}}
\newcommand{\R}{\mathbb{R}}
\newcommand{\C}{\mathbb{C}}
\newcommand{\F}{\mathbb{F}}
\newcommand{\A}{\mathbb{A}}
\newcommand{\PP}{\mathbb{P}}
\newcommand{\E}{\mathbb{E}}
\newcommand{\eps}{\varepsilon}
\newcommand{\dd}{\,\mathrm{d}}
\begin{document}
\section*{UnsolvedMath ID: 30002788}
\textbf{Source problem code:} catalogue entry
\subsection{Definitions and mathematical statement}
The supplied question is: ``Is the sum of two linearly computable forms always linearly computable?''

A \emph{form} ordinarily means a homogeneous polynomial. However, the supplied record does not define \emph{linearly computable}, specify the coefficient field or characteristic, or state the degrees and variable spaces of the two forms. Consequently, it does not determine a complete mathematical formulation.

The closure question can be recorded schematically as follows. Let \(\mathcal H\) denote the intended class of forms for which addition is defined, and let \(\mathcal L\subseteq\mathcal H\) denote those forms satisfying the intended, unspecified definition of linear computability. Is it true that
\[
\forall F,G\in\mathcal L,\qquad F+G\in\mathcal L?
\]
Here \(\mathcal H\) and \(\mathcal L\) are placeholders for missing source definitions, not independently specified mathematical objects. This display encodes only the wording of the question.

A precise statement requires the definition of linear computability and its ambient hypotheses. In particular, the source must clarify whether the summands have the same degree, whether they use a common variable space or disjoint sets of variables, and whether the zero form is included when cancellation occurs. No choice among these interpretations is made here.

The supplied curation labels the problem \emph{partially solved}. Its assessment says that the record itself gives a counterexample to the claim that all forms are linearly computable, that additivity and closure are established only in special cases, and that the stated closure question remains open. The supplied material includes neither that counterexample nor the definitions and hypotheses of those special cases. These are supplied literature claims, not mathematical facts independently verified in this draft.
\subsection{Short English statement}
The question asks whether adding two linearly computable forms always produces another linearly computable form. The supplied record leaves the key property and the permitted summands undefined, so only a schematic closure statement can be given. The supplied partially solved status and claim that the closure question remains open require independent review.
\subsection{Sources}
\begin{itemize}
\item Supplied local UnsolvedMath registry record, permanent ID \(30002788\), source problem code \texttt{catalogue entry}, in \texttt{data/unsolvedmath/problems.json}; collection: Oberwolfach Reports Open Problems. The record gives proposed year 2015 and supplies no named proposer. The local registry was not independently inspected for this draft.
\item Supplied source citation: \emph{Mini-Workshop: Ideals of Linear Subspaces, Their Symbolic Powers and Waring Problems} (2015), Oberwolfach Reports, pp.~489--532, DOI \texttt{10.4171/owr/2015/9}. \url{https://doi.org/10.4171/owr/2015/9}. The supplied background describes Oberwolfach Reports as documenting mathematical workshops and their open-problem discussions; no problem-specific page or passage is supplied.
\item Supplied literature source, labeled ``Teitler, sufficient conditions for Strassen additivity'': \url{https://arxiv.org/abs/1604.07691}. No theorem number, exact hypotheses, or supporting excerpt is supplied.
\item Supplied literature source, labeled ``original OWR report'': \url{https://doi.org/10.4171/owr/2015/9}.
\item Literature assessment date: \texttt{2026-08-21}. This is supplied dated curation, described in the record as literature triage derived from the supplied OWR curation, not an independently verified assessment. It labels the problem partially solved, reports a counterexample to universal linear computability and special cases of additivity and closure, and describes the stated closure question as remaining open. The record expressly calls for independent verification of time-sensitive status and cited claims.
\end{itemize}
\end{document}
